\clearpage\pagestyle{otMain}\setlength{\otBlocEspacement}{2.5mm}\fontsize{11}{15}\selectfont\corrigesChapitre{12}{Repérage et grandeurs en fonction d’une autre}
\begin{blocLivre}{cor-12-1}\begin{corrige}{1}{Lire deux points}\demarche{Projeter sur chacun des axes.}
On lit d’abord l’abscisse puis l’ordonnée : $\res{A(2;3)}$ et $\res{B(-3;1)}$.
\label{fig-cor-12-1}\figureGeo[corrige]{rep-1}
\end{corrige}\end{blocLivre}
\begin{blocLivre}{cor-12-2}\begin{corrige}{2}{Placer un point}\demarche{Suivre l’ordre abscisse puis ordonnée.}
Aller deux unités à gauche puis trois vers le bas. Le point est dans la zone où les deux coordonnées sont négatives.
\label{fig-cor-12-2}\figureGeo[corrige]{rep-2}
\end{corrige}\end{blocLivre}
\begin{blocLivre}{cor-12-3}\begin{corrige}{3}{Sur les axes}\demarche{Utiliser le rôle du zéro.}
D est sur l’axe vertical, E sur l’axe horizontal. Une coordonnée nulle supprime le déplacement dans la direction correspondante.
\label{fig-cor-12-3}\figureGeo[corrige]{rep-3}
\end{corrige}\end{blocLivre}
\begin{blocLivre}{cor-12-4}\begin{corrige}{4}{Inverser les coordonnées}\demarche{Ne pas permuter les axes.}
Leurs abscisses et ordonnées diffèrent ; les deux points sont $\res{\text{distincts}}$. Les coordonnées se lisent dans un ordre fixé.
\label{fig-cor-12-4}\figureGeo[corrige]{rep-4}
\end{corrige}\end{blocLivre}
\begin{blocLivre}{cor-12-5}\begin{corrige}{5}{Lire une valeur de graphique}\demarche{Lire l’ordonnée associée.}
À l’abscisse 2, on rejoint la courbe puis l’axe vertical : $\res{10\,^\circ\mathrm C}$.
\label{fig-cor-12-5}\figureGeo[corrige]{rep-5}
\end{corrige}\end{blocLivre}
\begin{blocLivre}{cor-12-6}\begin{corrige}{6}{Construire un tableau}\demarche{Substituer chaque entrée.}
On obtient $\res{4,\ 6,\ 10}$ euros ; les points sont $(0;4)$, $(1;6)$ et $(3;10)$. Les valeurs calculées sont ajoutées dans le repère.
\label{fig-cor-12-6}\figureGeo[corrige]{rep-6}
\end{corrige}\end{blocLivre}
\begin{blocLivre}{cor-12-7}\begin{corrige}{7}{Une courbe de carré}\demarche{Calculer avant de tracer.}
Les aires valent $\res{0,1,4,9}$ cm$^2$. Les points $(0;0)$, $(1;1)$, $(2;4)$ et $(3;9)$ ne sont pas alignés : l’aire n’est pas proportionnelle au côté.
\label{fig-cor-12-7}\figureGeo[corrige]{rep-7}
\end{corrige}\end{blocLivre}
\begin{blocLivre}{cor-12-8}\begin{corrige}{8}{Coordonnées décimales}\demarche{Adapter les déplacements au pas du quadrillage.}
Chaque petit intervalle vaut 0,5. H est à trois pas à gauche et cinq vers le haut ; I à cinq pas à droite et un vers le bas.
\label{fig-cor-12-8}\figureGeo[corrige]{rep-8}
\end{corrige}\end{blocLivre}
\begin{blocLivre}{cor-12-9}\begin{corrige}{9}{Lire plusieurs moments}\demarche{Distinguer hausse, palier et baisse.}
On lit $\res{10,\ 30,\ 20\ \mathrm L}$. Le segment horizontal de 2 à 4 minutes indique un volume constant de 30 L pendant cette période.
\label{fig-cor-12-9}\figureGeo[corrige]{rep-9}
\end{corrige}\end{blocLivre}
\begin{blocLivre}{cor-12-10}\begin{corrige}{10}{Des points alignés}\demarche{Contrôler le passage par l’origine.}
La droite passe par $(0;3)$, et non par l’origine. La relation est $\res{\text{non proportionnelle}}$ : pour une entrée nulle, la sortie vaut 3.
\label{fig-cor-12-10}\figureGeo[corrige]{rep-10}
\end{corrige}\end{blocLivre}
\begin{blocLivre}{cor-12-11}\begin{corrige}{11}{Un rectangle dans le plan}\demarche{Croiser les deux coordonnées utiles.}
Les côtés suivent les axes. D a l’abscisse de A et l’ordonnée de C : $\res{D(-3;3)}$. Les côtés opposés sont parallèles.
\label{fig-cor-12-11}\figureGeo[corrige]{rep-11}
\end{corrige}\end{blocLivre}
\begin{blocLivre}{cor-12-12}\begin{corrige}{12}{Retrouver une formule}\demarche{Relier la formule à une lecture.}
La formule est $\res{2t+5}$. Pour 4 h, $2\times4+5=\res{13}$ euros. Le graphique donne le même résultat.
\label{fig-cor-12-12}\figureGeo[corrige]{rep-12}
\end{corrige}\end{blocLivre}
\begin{blocLivre}{cor-12-13}\begin{corrige}{13}{Comparer deux tarifs}\demarche{Comparer les sorties pour une même entrée.}
A donne $0,3,6,9,12,15$ ; B donne $6,7,8,9,10,11$. Les prix sont égaux pour $\res{3\ \mathrm h}$ : 9 euros. Pour 2 h, A est moins cher ; pour 4 h, B l’est.
\label{fig-cor-12-13}\figureGeo[corrige]{rep-13}
\end{corrige}\end{blocLivre}
\begin{blocLivre}{cor-12-14}\begin{corrige}{14}{Lire une tendance sans extrapoler}\demarche{Séparer observations et modèle supposé.}
Les points connus sont $(1;2)$, $(2;5)$, $(3;5)$, $(4;8)$. Ils ne sont pas alignés avec l’origine : ce n’est pas une proportionnalité. Sans règle ou mesure supplémentaire, la valeur à 5 minutes est $\res{\text{indéterminée}}$.
\label{fig-cor-12-14}\figureGeo[corrige]{rep-14}
\end{corrige}\end{blocLivre}
\begin{blocLivre}{cor-12-15}\begin{corrige}{15}{Une réserve qui baisse}\demarche{Préciser le domaine où le modèle décrit la situation.}
Les volumes sont $\res{24,18,12,6,0}$ L. La formule est $\res{24-3t}$ pour $0\le t\le8$. Après huit minutes, le réservoir est vide : prolonger la droite donnerait des volumes négatifs sans sens pour cette situation.
\label{fig-cor-12-15}\figureGeo[corrige]{rep-15}
\end{corrige}\end{blocLivre}
